% ADDED figure (2026-11-02): the diagrams of one example of §C6b.11; generated by qft-C/tools/c6b11/figs/make_figs.sh
\tikzset{
  fdl/.style={line width=0.8pt, black},
  fdx/.style={circle, fill=black, inner sep=0pt, minimum size=3.6pt},
  fdv/.style={circle, fill=figB, draw=figB, inner sep=0pt, minimum size=5pt},
  fdlab/.style={font=\scriptsize, inner sep=2.6pt},
  fdmom/.style={->, >=Stealth, figO, line width=0.6pt},
}
\tikzset{
  fda/.style={fdl, postaction={decorate}, decoration={markings, mark=at position 0.55 with {\arrow{Stealth[length=5pt]}}}},
  fdS/.style={font=\small},
}
\begin{tikzpicture}[font=\small]
  \draw[fdl] (-1.2,0.8)--(0,0)--(1.2,0.8); \draw[fdl] (-1.2,-0.8)--(0,0)--(1.2,-0.8);
  \node[fdv] at (0,0){};
  \node[fdlab,left] at (-1.2,0.8){$p_A$}; \node[fdlab,left] at (-1.2,-0.8){$p_B$};
  \node[fdlab,right] at (1.2,0.8){$p_1$}; \node[fdlab,right] at (1.2,-0.8){$p_2$};
  \node[fdS] at (0,-1.45) {$\phi\phi\to\phi\phi$: $i\mathcal M=-i\lambda$, $S=1$};
\end{tikzpicture}
